% from ega4-errata-addenda-list3-fr.tex
\oldpage[IV]{353}
que la restriction $f|V:V\to Y$ est une immersion ouverte; donc il y a
un voisinage ouvert $U$ de $y\in Y$ tel que $f|V$ soit un isomorphisme
de $V$ sur le sous-préschéma induit par $Y$ sur $U$; l'isomorphisme
réciproque $g$ est alors une $U$-section de $f^{-1}(U)$ telle que
$g(y)=x$.

\textup{II)} L'obtention d'une $U$-section $g$ telle que $g(y)=x$ étant
un problème local sur $X$ et sur $Y$, on peut supposer que $X$ et $Y$
sont affines et que $f$ est de type fini, donc que $X$ est un
sous-préschéma fermé d'un $Y$-préschéma $Z$ de présentation finie sur
$Y$ (par exemple de la forme $Y[T_1,\ldots,T_n]$). Posons
\[
 Y'=\operatorname{Spec}(\sh{O}_{Y,y}),\qquad
 X'=X_{(Y')},\qquad Z'=Z_{(Y')},
\]
de sorte que $X'$ est un sous-préschéma fermé de $Z'$. En outre, si
$f'=f_{(Y')}$, on a $f'^{-1}(y)=f^{-1}(y)$ (I, 3.6.5) et $x$ est donc
isolé dans $f'^{-1}(y)$. Puisque $Y'$ est normal, on peut appliquer le
résultat de I) à $f'$, et puisque $Y'$ est le seul voisinage de $y$ dans
$Y'$, il existe une $Y'$-section $g'$ de $X'$ telle que $g'(y)=x$, que
l'on peut évidemment aussi considérer comme une $Y'$-section de $Z'$.
Puisque le morphisme structural $h:Z\to Y$ est de présentation finie,
on peut appliquer ({\hyperref[canon-index.ega4.c04.thm.8-8-2]{8.8.2}}, (i) et
(ii)) suivant la méthode de
({\hyperref[canon-index.ega4.c04.statement.8-1-2]{8.1.2}}, a)), d'où l'on conclut
à l'existence d'un voisinage ouvert $U$ de $y$ dans $Y$ et d'une
$U$-section $g$ de $h^{-1}(U)$ telle que $g(y)=x$. Il reste à voir
qu'en fait $g$ est une $U$-section de $f^{-1}(U)$. Mais puisque
$f^{-1}(U)$ est un sous-préschéma fermé de $h^{-1}(U)$ et que $U$ est
réduit, il suffit pour cela (I, 5.2.2) de voir que
$g(U)\subset f^{-1}(U)$. Or, puisque $U$ est intègre et $f^{-1}(U)$
fermé dans $h^{-1}(U)$, il suffit de prouver que si $\eta$ est le point
générique de $Y$ (et de $U$), on a $g(\eta)\in f^{-1}(U)$. Or, cela
résulte de ce que la restriction de $g$ à $Y'$ est une $Y'$-section de
$X'\subset f^{-1}(U)$ et de ce que $\eta\in Y'$.
\end{proof}

\paragraph{(Err\textsubscript{IV}, 31).}
Dans ({\hyperref[canon-index.ega4.c04.subsec.10-4]{IV, 10.4}}), ajouter, après la
proposition {\hyperref[canon-index.ega4.c04.statement.10-4-11]{(10.4.11)}} :

\begin{corollary}[10.4.12]
\EGAReaderTarget{ega4.c04.cor.10-4-12}{IV.10.4.12}
Soient $S$ un préschéma, $X$ un $S$-préschéma de présentation finie,
$f$ un $S$-endomorphisme radiciel de $X$.
\begin{enumerate}[label=\textup{(\roman*)}]
\item L'endomorphisme $f$ est bijectif et fini (donc un homéomorphisme
universel {\hyperref[canon-index.ega4.c04.prop.2-4-5]{(2.4.5)}}).
\item Supposons de plus que $X$ soit réduit, que les anneaux locaux de
$S$ soient noethériens et universellement japonais, et que $f$ soit
birationnel (condition automatiquement vérifiée si pour tout point
maximal $x$ de $X$, $k(x)$ est un corps parfait). Alors $f$ est un
$S$-automorphisme.
\end{enumerate}
\end{corollary}

\begin{proof}
\textup{(i)} La question étant locale sur $S$, on peut se borner au cas
où $S=\operatorname{Spec}(A)$ est affine. Il existe alors une
sous-$\mathbb Z$-algèbre de type fini $A_0$ de $A$ telle qu'en posant
$S_0=\operatorname{Spec}(A_0)$, on ait
$X=X_0\times_{S_0}S$, où $X_0$ est un préschéma de type fini sur $S_0$,
et que $f$ se déduise d'un $S_0$-endomorphisme $f_0$ de $X_0$ par le
changement de base $S\to S_0$
{\hyperref[canon-index.ega4.c04.prop.8-9-1]{(8.9.1)}}; on peut en outre supposer
$f_0$ radiciel ({\hyperref[canon-index.ega4.c04.thm.8-10-5]{8.10.5}}, (vii)). Il
suffira de prouver que $f_0$ est bijectif et fini, et on peut donc
supposer désormais que $A$ est une $\mathbb Z$-algèbre de type fini, ce
qui entraîne que $S$ et $X$ sont des préschémas excellents
({\hyperref[canon-index.ega4.c04.statement.7-8-3]{(7.8.3)}} et
{\hyperref[canon-index.ega4.c04.prop.7-8-6]{(7.8.6)}}) et en particulier que $X$
est un préschéma universellement japonais
({\hyperref[canon-index.ega4.c04.cor.5-11-4]{(5.11.4)}} et
({\hyperref[canon-index.ega4.c04.statement.7-8-3]{7.8.3}}, (vi))). On peut en outre
supposer $X$ réduit; en effet,
$f_{\mathrm{red}}:X_{\mathrm{red}}\to X_{\mathrm{red}}$ est un
$S_{\mathrm{red}}$-endomorphisme radiciel (I, 5.1.6); si l'on prouve que
$f_{\mathrm{red}}$ est fini, il en sera de même de $f$ : en effet,
puisque $f$ est radiciel, il est séparé
{\hyperref[canon-index.ega4.c04.lem.1-8-7-1]{(1.8.7.1)}}; comme
$f_{\mathrm{red}}$ est propre il en est de même de $f$ (II, 5.4.6);
mais puisque $f$ est évidemment quasi-fini, il est fini (III, 4.4.2).

